% MCDS Capstone Spring Report Document
% Kyle Dotterrer
% 26 April, 2021

%% ==================================================================
%% BOILERPLATE
%% ==================================================================

\documentclass{vldb}

\newcommand{\paperTitle}{JIT Compilation for User-Defined Functions\\ 
    \large Capstone Requirements Document\\
    \today}
\newcommand{\paperKeywords}{}
\newcommand{\paperAuthors}{}

\setlength{\paperheight}{11in}
\setlength{\paperwidth}{8.5in}

\usepackage[
            bookmarks=true,
            bookmarksopen=true,
            pdfhighlight=/I,
            pdfpagemode=UseOutlines,
            linkcolor=blue,
            pdfborder={ 0 0 0 },
            pageanchor=false]{hyperref}
\hypersetup{
    pdfauthor = {\paperAuthors},
    pdftitle = {\paperTitle},
    pdfkeywords = {\paperKeywords},
    pdfborder={ 0 0 0 }
}

\usepackage{amsmath}
\usepackage{amssymb}
\usepackage{times}
\usepackage{boxedminipage}
\usepackage{xspace}
\usepackage{array}
\usepackage{epsfig}
\usepackage{calc}
\usepackage{multirow}
\usepackage{rotating}
\usepackage{enumitem}
\usepackage[usenames,dvipsnames]{color}
\usepackage{tabularx}
\usepackage{tikz}
\usepackage{wasysym }
\usepackage[hyphenbreaks]{breakurl}
\usepackage{xcolor}
\usepackage[font={small}]{caption}
\usepackage{alltt}
\usepackage{setspace}
\usepackage{wrapfig}
\usepackage{etoolbox}

% Font selection
\usepackage[final]{microtype}
\usepackage[scaled]{inconsolata}
\usepackage[T1]{fontenc}
\usepackage{graphicx}
\usepackage{soul}
\usepackage{pifont}

% Other misc. things
\usepackage{float}
\usepackage{bookmark}
\usepackage{amsmath,amsfonts,amssymb}
\usepackage{latexsym}
\usepackage{multirow}
\usepackage{emptypage}
\usepackage{booktabs}
\usepackage{enumerate}
\usepackage{braket}
\usepackage{balance}

%% ==================================================================
%% MACROS
%% ==================================================================

\newcommand{\subparagraph}{}

\newcommand{\sysname}{NoisePage\xspace}

% Framework Components
\newcommand{\dbTranslator}{Translator\xspace}
\newcommand{\dbCompiler}{Compiler\xspace}
\newcommand{\dbPlanner}{planner\xspace}
\newcommand{\dbInterpreter}{interpreter\xspace}
\newcommand{\dbAnalyzer}{analyzer\xspace}
\newcommand{\dbPCQ}{PCQ\xspace}
\newcommand{\dbAQP}{AQP\xspace}
\newcommand{\dbHook}{hook\xspace}
\newcommand{\dbHooks}{hooks\xspace}

\newcommand{\dbSQL}[1]{\texttt{\textbf{#1}}\xspace}
\newcommand{\dbTable}[1]{\texttt{\MakeUppercase{#1}}\xspace}
\newcommand{\dbColumn}[1]{\texttt{col#1}\xspace}
\newcommand{\dbRawColumn}[1]{\texttt{#1}\xspace}
\newcommand{\dbConfig}[1]{\textsf{#1}\xspace}

\newcommand{\dbmsOracle}{Oracle\xspace}
\newcommand{\dbmsDBTwo}{DB2\xspace}
\newcommand{\dbmsTeradata}{Teradata\xspace}
\newcommand{\dbmsHANA}{HANA\xspace}
\newcommand{\dbmsMSSQL}{SQL Server\xspace}
\newcommand{\dbmsMySQL}{MySQL\xspace}
\newcommand{\dbmsPostgres}{PostgreSQL\xspace}
\newcommand{\dbmsMemSQL}{MemSQL\xspace}
\newcommand{\dbmsHyper}{HyPer\xspace}
\newcommand{\dbmsVector}{Vector\xspace}

\newcommand{\tplType}[1]{\texttt{#1}\xspace}
\newcommand{\tpl}{\texttt{TPL}\xspace}
\newcommand{\tbc}{\texttt{TBC}\xspace}
\newcommand{\dbCode}[1]{{\sffamily\small \textbf{#1}}\xspace}

\newcommand{\tidlist}{\texttt{TupleIdList}\xspace}
\newcommand{\tidlists}{\texttt{TupleIdLists}\xspace}

%% ==================================================================
%% PAPER CONTENT
%% ==================================================================

\begin{document}

\title{\paperTitle}

\numberofauthors{2}
\author{
    \alignauthor Kyle Dotterrer, Andrew Pavlo\\
    \affaddr{Carnegie Mellon University}\\
    \email 
        \{\href{mailto:kdottere@cs.cmu.edu}{kdottere},
        \href{mailto:pavlo@cs.cmu.edu}{pavlo}\}@cs.cmu.edu\\
}

% Remove the space for the copyright box for VLDB
\makeatletter
\def\@copyrightspace{\relax}
\makeatother

\maketitle

%% ==================================================================
%% ABSTRACT
%% ==================================================================

\section*{Abstract}

User-defined functions (UDFs) are a powerful tool for encapsulating complex application logic that may subsequently be embedded into SQL queries. Although they offer a practical complement to SQL, their use in production database environments is often limited because of their associated performance costs \cite{froid, apfel}. In this project, we evaluate the efficacy of methods intended to address the performance limitations of UDFs in the context of main-memory database management systems (DBMSs). We compare the performance of an existing, state-of-the-art technique, UDF inlining, with our own novel strategy that executes UDFs as native code via JIT-compilation. Our evaluation efforts are ongoing.

%% ==================================================================
%% INTRODUCTION
%% ==================================================================

\section{Introduction}

TODO

The remainder of this document is structured as follows. Section 2 begins with a discussion of our project goals and hypotheses. Next, in Section 3, we provide an overview of relevant prior work and its relationship with our project. Section 4 outlines our project goals for the Spring term and discusses the degree to which we have met them. Sections 5 and 6 provide overviews of our system and experimental designs, respectively. In sections 7 and 8 we present the results and error analysis for the results we have achieved thus far. Section 9 discusses the entirety of our design, implementation, and evaluation process, to the extent currently possible. In Section 10 we present some of the lessons we have learned thus far in the course of our project. Section 11 presents our current plan for the work that remains to be done in our project. Finally, we conclude in Section 12.

%% ==================================================================
%% PROJECT GOAL
%% ==================================================================

\section{Project Goal and Hypotheses}

Our goal in this project is to evaluate methods for improving UDF execution performance in database systems. We compare the performance of the existing state-of-the-art (\cite{froid, apfel, aggify}) with our own novel implementation of UDF JIT-compilation.

We evaluate the following hypotheses through our implementation efforts and subsequent evaluation:

\begin{itemize}
    \item Integrating JIT compilation of UDFs into an existing DBMS with a mature execution engine based on both interpretation and native execution is feasible; prior work by Nayak \cite{tanuj} provides strong evidence that this is the case. In this context, success will be measured by the ability to merge the implementation of UDF JIT compilation with future releases of NoisePage, observing the existing robustness and coding standards of the project.
    \item JIT compilation of UDFs to native machine code will yield performance that surpasses that achieved by traditional approaches such as UDF interpretation. Furthermore, the performance of UDF compilation will match or exceed that demonstrated by alternative, state-of-the-art approaches such as UDF inlining \cite{apfel, froid, aggify}. In this context, success may be measured quantitatively: query execution time will be the primary metric of interest for analytical workloads while throughput (transactions per second) will be the analog for transactional workloads.
\end{itemize}

%% ==================================================================
%% RELATIONSHIP TO PRIOR WORK
%% ==================================================================

\section{Relationship to Prior Work}

Our goal in this project is to compare various UDF execution methods, namely JIT compilation and UDF inlining. Our comparison's primary target is system performance, which we will measure by the metrics appropriate for specific workload characteristics. Our metric of interest will be throughput for transactional workloads, while for analytical workloads, our metric of interest will be total query processing time (latency). In the following subsections, we survey the prior work in two relevant areas of ongoing research. First, we provide an introduction to query compilation techniques in database systems. We follow this with a discussion of existing methods for improving UDF execution performance via inlining.

\subsection{Query Compilation}

In traditional disk-oriented systems, the DBMS implements query execution entirely in software. The canonical execution strategy employed by many database systems begins as the optimizer converts the abstract syntax tree produced by parsing the query text into a physical plan tree. Each node in this plan tree represents a physical operator that implements one or more logical operations specified in the input query, such as a \dbCode{SELECT} operation, application of a filter (\dbCode{WHERE}), or a \dbCode{JOIN}. With the plan tree constructed, the DBMS executes the query by evaluating each operator in the physical plan \cite{graefe93, graefe94, dbconcepts}. 

In this strategy, the DBMS implements its execution engine in terms of high-level, general-purpose code. This degree of generality is excellent from a software engineering perspective because it makes query execution easy to reason about. However, this comprehensibility comes at the cost of runtime performance. The generality of the code in this execution technique imposes runtime overhead by way of the many levels of indirection that it introduces, resulting in a large number of native instructions being executed \cite{freedman14}. The number of executed instructions is often not a limiting factor in disk-oriented systems' performance because disk IO operations dominate the total query processing time. For in-memory systems where persistent storage no longer imposes this bottleneck, several strategies are available to reduce the number of instructions executed in query processing.

Query compilation is one such query execution strategy that is growing in popularity. At a high level, query compilation aims to reduce the number of instructions executed by generating specialized code for each query that the system encounters. 

There are many approaches to implementing query compilation, and the techniques for doing so have evolved rapidly in the last decade. One of the earliest implementations of query compilation in the modern era was accomplished by the HIQUE system \cite{systemr}. In HIQUE, the execution engine first translates an input query to C++ code. It subsequently compiles this C++ code to native code by spawning an external process (GCC \cite{gcc}) \cite{krikellas10}. It is important to note here, however, that HIQUE does not implement translation for arbitrary SQL queries but instead relies on pre-defined C++ templates that it applies to a limited subset of input queries.

More advanced systems eschew using an external compiler to avoid the high costs associated with compilation under this model. Spawning an external compilation process introduces the latency associated with operating system process management into the critical path of query execution. Furthermore, utilizing an external compiler like GCC to compile specialized query code discards much of the information that the DBMS has already `learned' about the query while constructing its query plan and forces the compiler's frontend to repeat much of this work. The HyPer main memory DBMS \cite{hyper} is one example of a system that avoids these issues by implementing query compilation without using an external compiler. HyPer generates LLVM \cite{LLVM} intermediate representation from the input query and subsequently utilizes the existing LLVM infrastructure to JIT-compile the IR to native machine code that implements the query \cite{neumann11, kemper11}.

The SingleStore \cite{singlestore} execution engine represents a more advanced approach to compilation. It first translates SQL queries to an internal domain-specific language (DSL) that then undergoes a second stage of translation to LLVM intermediate representation. The execution engine then JIT-compiles this intermediate representation to native code \cite{paroski16}. This additional step in the query compilation process facilitates greater developer productivity by making it easier to reason about the code generated by the execution engine's upper layers.

Carnegie Mellon Database Group's NoisePage system \cite{noisepage} implements an approach to query compilation that is similar to SingleStore's. The first step in NoisePage's query processing pipeline is translation of the input query to an internal domain-specific language. The execution engine then compiles this DSL to a lower-level bytecode representation. Once this bytecode is generated, the NoisePage execution engine can either run the bytecode directly on an interpreter or translate this bytecode to LLVM IR that is then JIT-compiled to efficient machine code \cite{pcq}.

\subsection{UDF Inlining}

Despite recent advances in compilation techniques, the most popular current approach to addressing the performance limitations of user-defined functions does not involve native compilation. Instead, most of the attention in this space is currently devoted to the technique of UDF inlining. In general, inlining entails translating a UDF into a relational algebraic expression that can then be embedded directly into the calling SQL query. Inlining a UDF allows the DBMS's existing query optimization component to optimize the logic of an imperative UDF in the same way it optimizes declarative SQL.

Ramachandra et al. implement this technique's archetype in their Froid system within Microsoft SQL Server \cite{sqlserver}. Froid executes UDF inlining by first decomposing an individual UDF into distinct \textit{regions} that it then translates into equivalent relational expressions. The Froid framework then relies on SQL Server's \dbCode{APPLY} operator to compose these relational expressions into a single relational expression that encodes the entirety of the UDF's logic \cite{froid}. Gupta et al. extend this work on Froid with their Aggify framework to add support for cursor loops. Aggify realizes its support for cursor loops by translating them to equivalent custom aggregations \cite{aggify}. 

Duta et al. pursue a distinct approach to UDF inlining in their Apfel framework \cite{apfel}. Where Ramachandra et al. implement Froid as an integrated component of the Microsoft SQL Server DBMS, Duta et al. design their Apfel framework as a standalone system that compiles imperative UDFs written in PL/pgSQL \cite{plpgsql} to standard SQL. Any database system may then process the resulting SQL query in the usual manner. The Apfel framework's defining feature is its use of recursive common table expressions (SQL's \dbCode{WITH RECURSIVE} syntax) to reproduce arbitrary imperative loop logic in SQL.

%% ==================================================================
%% SPRING DEVELOPMENT GOALS
%% ==================================================================

\section{Spring Development Goals}

At the outset of this project, we set the goal to complete implementation of our UDF JIT compilation technique by the conclusion of the Spring semester. This goal was not amended at the semester midterm because we observed that we were progressing at approximately the expected rate. Now, at the conclusion of the semester, we can happily report that we did meet our stated goal: we successfully merged support for UDF JIT compilation into the NoisePage DBMS project codebase, implying that the system now supports our novel UDF execution strategy. Obviously, this constitutes just one step towards our ultimate goal of an experimental evaluation, but we believe that this represents the majority of the engineering effort presented by the project in its totality.

%% ==================================================================
%% SYSTEM DESIGN
%% ==================================================================

\section{System Design}

To perform a robust evaluation of UDF execution strategies, we require the existence of a system that implements both techniques under test - UDF inlining and UDF JIT compilation. To meet this requirement, we implement both of these techniques in the NoisePage DBMS. One system feature that is common to both UDF inlining and UDF JIT compilation is the existence of an \textit{Execution engine} with support for compiled query execution. In the subsections that follow, we provide details of NoisePage's \textit{Execution Engine} subsystem, followed by the specifics of how inlining and JIT compilation interact with this subsystem to execute user-defined functions.

\subsection{Compiled Query Execution}

Carnegie Mellon Database Group's NoisePage system implements compiled query execution via a data-centric code-generation engine. The query compilation process accepts as input a physical query plan (produced by the \textit{Optimizer}) and is responsible for implementing the logic that executes it. The high-level steps that take place during query compilation and execution with the NoisePage \textit{Execution Engine} are as follows:

\begin{itemize}
    \item Translate the input physical query plan to an internal domain specific language, TPL \cite{pcq}. This domain-specific language resembles the C programming language in its syntax, but introduces many database-specific features such as native support for database operators and parallelization. It is also during this stage of query processing that the operators in the physical query plan are grouped into \textit{pipelines} \cite{rof}. Each pipeline is further partitioned into \textit{functions} that implement various aspects of the pipeline, such as setup, teardown, and serial or parallel execution. The \textit{functions} form the unit of work for later stages of the compilation process as well as the atomic unit of execution.
    \item Once the physical query plan has been translated to TPL, the \textit{Execution Engine} performs a further translation step, lowering the TPL DSL to a domain-specific bytecode, TBC \cite{pcq}. The translation from TPL to TBC resembles the compilation of a source language like C to a machine assembly language. We perform transformations such as rewriting loop constructs in terms of conditional tests and branch instructions, and insert bytecode instructions for explicit management of automatic (local) variables.
    \item Now that query is expressed in TBC, the NoisePage \textit{Execution Engine} may proceed in one of two ways, depending on the current configuration of the system. The first option is to \textit{Interpret} the bytecode directly on a custom interpreter. In this execution mode, the TBC is never lowered to native machine code, and instead the remaining stages of query execution take place entirely in software. The other option is to continue the lowering process and ultimately \textit{Compile} the query to native machine code via the LLVM JIT infrastructure. This latter option is the path of interest for the purposes of our project, so we will assume that the system proceeds in this configuration for the remainder of this discussion.
    \item The NoisePage \textit{Execution Engine} does not translate directly from the bytecode representation of the query to native machine code, and instead relies on the LLVM compiler infrastructure to assist in this process. For this reason, the next step in query compilation is translation of the TBC bytecode to the LLVM intermediate representation (IR).
    \item Now that the input query is expressed in LLVM IR, the NoisePage \textit{Execution Engine} utilizes the LLVM JIT API to compile each function in the query to native machine code, completing the compilation process.
    \item Once each \textit{function} is compiled to native machine code, the \textit{Execution Engine} executes the query by invoking each function in the proper order. In this context, function invocation refers to native function invocation, implying that the overhead of each invocation is minimal.
\end{itemize}

This completes the enumeration of the stages of query compilation in the NoisePage \textit{Execution Engine}. 

\subsection{UDF Execution}

Next we turn our attention to the specifics of how we implement UDF JIT compilation feature on top of this subsystem.

\subsection{UDF Inlining}

Finally, in this subsection we briefly describe how we implement UDF inlining on top of the existing query compilation infrastructure in NoisePage. It should be noted that system design work for this feature is still in its early phases, and that the following description may be subject to substantial revision.

%% ==================================================================
%% EXPERIMENTAL DESIGN
%% ==================================================================

\section{Experimental Design}

In this section we discuss relevant aspects of our experimental design. As mentioned in previous sections, the majority of our time this semester was spent on design and implementation work for UDF JIT compilation, rather than the conduct of our evaluation. Accordingly, the setup described below merely represents the proposed design of our evaluation, rather than the setup which we have actually used. See the \textit{Future Work} section for further details regarding our plans for the remainder of the project.

We partition our experimental evaluation into two components: targeted performance evaluation and representative performance evaluation. We describe the features of each of these components in the subsections below. Note that the ordering of the two parts of our experiment is immaterial - they are independent. Following this, we provide an overview of the metrics of interest in our experiment.

\subsection{Computational Overhead Evaluation}

In the first part of our evaluation, we will investigate the degree of computational overhead imposed by each of the UDF execution methods that we consider. This entails constructing targeted benchmarks that exercise one or more aspects of the imperative programming model provided by UDFs and observing the resulting performance distinctions. These benchmarks will consist primarily of constructs that are well-suited to imperative formulation. This aspect of our evaluation is important because it isolates the algorithmic costs associated with each execution technique and allows us to make inferences regarding their future potential for improvement.

\subsection{Production Workload Evaluation}

In the second component of our evaluation, we will investigate the performance of UDF execution techniques in the context of inputs that are representative of those found in production environments. This aspect of our evaluation is important because it answers questions regarding the current applicability of the techniques under evaluation.

The nature of this evaluation implies that we require the ability to identify the types of UDFs that are typically used in production environments. In practice, identifying the appropriate benchmarks to use for this part of the evaluation is not straightforward. For many database workloads, industry-standard benchmarks exist to measure the performance of a particular system. For instance, the Transaction Processing Council publishes both the TPC-C \cite{tpcc} and TPC-H \cite{tpch} benchmarks. These benchmarks are widely accepted within the database systems community as effective measures of online transaction processing and analytical workload performance, respectively. To our knowledge, no such standardized benchmarks exist that emphasize the use of user-defined functions. For this reason, prior work in this area has typically involved auxiliary efforts to construct benchmarks appropriate for user-defined functions. For instance, Ramachandra et al. analyze customer workloads from Microsoft's Azure SQL Database \cite{azure} and rank the workloads in descending order of UDF-prevalence in their evaluation of Froid \cite{froid}. Nayak pursues an open-source solution to the same problem by analyzing queries scraped from public GitHub \cite{github} repositories \cite{tanuj}. Our evaluation will likely involve a combination of both approaches.

\subsection{Evaluation Metrics}

We are interested in two classes of metrics in our evaluation of UDF execution techniques: performance improvements and resource consumption. These classes of metrics are in line with those evaluated in prior work on this topic \cite{froid, apfel}.

In the context of performance improvements, we are interested in the following metrics:

\begin{itemize}
    \item \textit{Compilation Time}: The total time taken to compile the query from a physical plan to native machine code. We hypothesis that the different techniques under evaluation introducing differing degrees of compile-time overhead to the system. This metric is measured in milliseconds.
    \item \textit{Execution Time}: The total time taken to execute the compiled query. Again, we hypothesize that the different techniques under evaluation will result in measurably different execution time even though both result solely in native code execution. This metric is measured in milliseconds.
    \item \textit{Combined Compilation and Execution Time}: The total time taken to compile and execute the query within the DBMS execution engine. Ultimately, it is the end-to-end time taken by the execution engine that determines the overall effectiveness of the technique under evaluation. This metric is measured in milliseconds.
\end{itemize}

In the context of resource consumption, we are interested in the following metrics:

\begin{itemize}
    \item \textit{CPU Utilization}: During compilation and execution, what percentage of available compute are we able to saturate? This metric will measured as a percentage.
    \item \textit{Memory Consumption}: During compilation and execution, what is the total volume of memory consumed by the DBMS execution engine? This metric is measured in megabytes (MB).
\end{itemize}

For both of these classes of metrics, we will evaluate each of the UDF execution techniques of interest on benchmarks designed according to the two classes described above.

%% ==================================================================
%% RESULTS
%% ==================================================================

\section{Results}

As discussed in previous sections, the entirety of our time this term was spent on design and implementation effort for our novel UDF execution technique - UDF JIT compilation. Accordingly, we have not yet conducted our proposed evaluation and therefore do not currently have any experimental results to present.

%% ==================================================================
%% ERROR ANALYSIS
%% ==================================================================

\section{Error Analysis}

As discussed in previous sections, the entirety of our time this term was spent on design and implementation effort for our novel UDF execution technique - UDF JIT compilation. Accordingly, we have not yet conducted our proposed evaluation and therefore do not currently have any error analysis to present.

%% ==================================================================
%% DISCUSSION
%% ==================================================================

\section{Discussion}

As discussed in previous sections, the entirety of our time this term was spent on design and implementation effort for our novel UDF execution technique - UDF JIT compilation. Accordingly, we have not yet conducted our proposed evaluation and therefore do not currently have any discussion of our experimental process, results, or error analysis to present.

%% ==================================================================
%% LESSONS LEARNED
%% ==================================================================

\section{Lessons Learned}

The following subsections summarize the lessons we learned during our system design and implementation work this semester.

\subsection{Working in a Mature Codebase}

Working in the context of a mature codebase presents challenges beyond those that might be expected for a standalone or "from scratch" systems design implementation project. The NoisePage project implements a fully-functional SQL-compliant in-memory database system, and accordingly it possesses many components that are distinct from yet dependent upon the query execution engine with which this project is concerned. Because of this, one must be careful when making any modification to the execution engine to ensure that the high-level semantics of the system are preserved. Additionally, verifying the correctness of the implementation through testing is made more difficult by the surrounding system; unless one thinks carefully about how each component of the new functionality interacts with the rest of the system and the interface it presents, it is altogether to easy to introduce hidden dependencies that make testing the new functionality in isolation impossible.

\subsection{Design versus Implementation}

This lesson perhaps should have been obvious from the outset, but it nevertheless proved to be something that we learned along the way. During the design phase of the project, it frequently felt like system design would represent the more arduous component of the first project task - once we had the design for UDF JIT compilation complete, implementation would be a simple matter of executing on our plan. Needless to say we were incorrect in this assumption. Implementation presented many subtle challenges that we failed to consider during our initial design process. Attempting to address such issues "on the fly" often led to low-quality solutions that introduced incidental complexity. For this reason, we often found it more productive to take additional time to deliberately revise our plan (with our hands off the keyboard) whenever such an issue arose.

%% ==================================================================
%% FUTURE WORK
%% ==================================================================

\section{Future Work}

Because our original project plan and timeline remain intact, our direction for future work is straightforward. Now that we have completed implementation of support for UDF JIT compilation, we turn our attention to implementation of UDF inlining support in the NoisePage DBMS. While this no doubt represents a non-trivial engineering effort, many of the system design issues in this space have been resolved by prior work (e.g. \cite{froid, apfel}). For this reason, we are confident that we can complete the implementation of UDF inlining support by the beginning of July, 2021, according to the schedule outline below.

\begin{center}
\begin{tabular}{ |c|l| } 
    \hline
    \textbf{Date} & \textbf{Milestone} \\ 
    \hline
    31 May 2021 & Begin system design \\
    \hline
    7 June 2021 & System design complete; integration plan \\
    \hline
    14 June 2021 & \textit{Binder} support implemented; tests passing \\
    \hline
    21 June 2021 & \textit{Optimizer} support implemented; tests passing \\
    \hline  
    28 June 2021 & Feature Complete \\
    \hline
\end{tabular}
\end{center}

With the completion of the system implementation work associated with UDF inlining, we will finally turn our attention to evaluation of UDF performance under each of the supported execution strategies. While this task constitutes the ultimate goal of the project, we hypothesize that it will require the less time and effort to complete than either of the two previous tasks. Our goal is to have results suitable for publication by the beginning of August, 2021. Therefore, we propose the following tentative timeline for the conduct of our evaluation.

\begin{center}
\begin{tabular}{ |c|l| } 
    \hline
    \textbf{Date} & \textbf{Milestone} \\ 
    \hline
    5 July 2021 & Complete experimental design for all cases \\
    \hline
    12 July 2021 &  All workload data collected \\
    \hline
    19 July 2021 & All benchmarks implemented \\
    \hline
    26 July 2021 & Benchmarks complete \\
    \hline
    2 August 2021 & Organization and Communication of Results \\
    \hline
\end{tabular}
\end{center}

Should we succeed in meeting the deadlines described above, we must address the question of what further work we plan to accomplish during the Fall semester. Fortunately, we foresee plenty of interesting work that remains to be done in the space of UDF execution strategies for database systems. Small, self-contained features related to UDF JIT compilation that might be of interest include the following:

\begin{itemize}
    \item Add support for dynamic execution of SQL statements
    \item Add support for additional PL/pgSQL constructs, such as \dbCode{RAISE} and \dbCode{EXCEPTION} semantics
    \item Add Support for vectorization / parallelization of TPL lambda execution
    \item Modify the implementation to unify the abstract syntax tree (AST) of the compiled UDF and the query in which it is embedded, potentially by precompiling the entire UDF to a dynamically-loaded LLVM function. Such a modification might enable significant performance gains as the LLVM optimizer would gain the ability to perform inter-procedural optimizations across the UDF call boundary.
\end{itemize}

Beyond these features, there are additional questions that we seek to address regarding the most effective way to integrate imperative programs with relational database systems in a more general sense. For instance, should our evaluation reveal that neither UDF JIT compilation nor UDF inlining represent a totally dominant execution strategy, is it possible for the DBMS query optimizer to select the more effective execution technique based solely on the information that is available at the time optimization is performed? We hypothesize that such a heterogenous execution strategy for UDFs will ultimately prove to offer the highest performance.

%% ==================================================================
%% CONCLUSION
%% ==================================================================

\section{Conclusion}

%% ==================================================================
%% ACKNOWLEDGEMENTS
%% ==================================================================

\section*{Acknowledgements}

I would like to thank my mentor Professor Andy Pavlo for his continued support and guidance throughout this project. Professor Pavlo was remarkably receptive to my proposed involvement in this project from the very beginning, and his passion for databases systems consistently serves as a bulwark against the discouragement posed by technical setbacks. In addition, I must acknowledge the prior work by former CMU student Tanuj Nayak, upon whose pioneering work with PL/pgSQL compilation this project is based. Finally, I would like to thank all of my colleagues in the CMU Database Group for their continued, groundbreaking work on the NoisePage DBMS and the inspiration this provides for my own development efforts.

%% ==================================================================
%% BIBLIOGRAPHY
%% ==================================================================

% NOTE: Had to add this mysterious command to compile...
\newcommand{\newblock}{}

\newpage
\balance
\bibliographystyle{abbrv}
\bibliography{report-spring}

\end{document}